\subsection{RESTORE Responsibilities}
RESTORE replaces each GSn whose corresponding (:ref:base.state.save_area.BASE_HEADER.GS_VALID:) bit is set and preserves each GSn whose bit is
clear. It validates every selected segment image before changing (:term:base.terms.architectural_state:).

RESTORE uses the state-block bitmap as authoritative when reconstructing extension state. It accepts only the
defined format and only bitmap bits that name CPUID-described components for that format. RESTORE raises
(:ref:base.events.EXCEPTION.INVALID_CONTROL_IMAGE:) before changing (:term:base.terms.architectural_state:) when the format is unrecognized or a set
bitmap bit does not name a CPUID-described component for that format. RESTORE of a clear, described component bit
installs that component\textquotesingle{}s initial state.

User-mode RESTORE ignores saved supervisor-only (:ref:base.state.save_area.BASE_HEADER.STATUS:) and (:ref:base.state.save_area.BASE_HEADER.DFA:) state. Supervisor RESTORE is an architectural
context transition: it restores (:ref:base.registers.STATE.STATUS.EDEPTH:), (:ref:base.registers.STATE.STATUS.UO:), and current (:ref:base.state.save_area.BASE_HEADER.DFA:) together with (:ref:base.state.save_area.BASE_HEADER.STATUS:) after validating the complete
event-state invariants. It requires saved (:ref:base.registers.STATE.STATUS.PM:)\texttt{=1}; (:ref:base.registers.STATE.STATUS.UO:) requires a valid live U bank, (:ref:base.state.save_area.BASE_HEADER.DFA:) requires an active event, and
inactive state requires (:ref:base.registers.STATE.STATUS.EDEPTH:)\texttt{=0}, (:ref:base.registers.STATE.STATUS.UO:)\texttt{=0}, and (:ref:base.state.save_area.BASE_HEADER.DFA:)\texttt{=0}. Ordinary WRSTATUS remains unable to change these (:term:base.terms.field|plural:).

The normative instruction entries are
(:ref:base.instructions.SAVE:) and
(:ref:base.instructions.RESTORE:).\par
